% =====================================================================
% MDBC --- Multi-Delta Blue Carbon: cross-delta transferability
% Target: Nature-portfolio style (Communications Earth & Environment).
% Standalone build (no publisher .cls required); retarget at submission.
% Build: pdflatex -> bibtex -> pdflatex x2
% =====================================================================
\documentclass[11pt]{article}

\usepackage[utf8]{inputenc}
\usepackage[T1]{fontenc}
\usepackage{lmodern}
\usepackage[a4paper,margin=1in]{geometry}
\usepackage{graphicx}
\graphicspath{{figures/}}
\usepackage{amsmath,amssymb}
\usepackage{booktabs}
\usepackage{siunitx}
\usepackage{authblk}
\usepackage{xcolor}
\usepackage[hidelinks]{hyperref}
\usepackage[round,authoryear]{natbib}
\usepackage{lineno}
\linenumbers

% editorial note macro (remove before submission)
\newcommand{\TODO}[1]{\textcolor{red}{[\textbf{TODO:} #1]}}
\newcommand{\PLACE}[1]{\textcolor{blue}{[\textit{#1}]}}

\title{\textbf{Cross-delta transferability limits of mangrove blue-carbon
stock models: a multi-delta benchmark with area-of-applicability and
conformal uncertainty}}

\author[1,2]{Md Rakib Hasan\thanks{Corresponding author:
\texttt{rakibhhridoy@fermium.systems}; ORCID: 0009-0002-4007-7590}}
\affil[1]{Fermium Systems, Dhaka 1207, Bangladesh}
\affil[2]{Department of Soil, Water and Environment, University of Dhaka,
Dhaka 1000, Bangladesh}

\date{\today}

\begin{document}
\maketitle

% ---------------------------------------------------------------------
\begin{abstract}
\noindent
Mangroves are among Earth's most carbon-dense ecosystems, and global maps of
their soil and biomass carbon now underpin carbon-crediting and conservation
decisions. Yet these maps are validated almost exclusively with random or
in-distribution cross-validation, which is inflated by spatial autocorrelation
and silent about whether a model trained on some regions can predict an
\emph{unsampled} delta. Here we build a multi-delta benchmark spanning
\PLACE{N} major mangrove systems across \PLACE{C} continents and evaluate
total ecosystem carbon (aboveground biomass carbon, AGB-C, plus soil organic
carbon, SOC) under three validation tiers --- random, spatial-block, and
leave-one-delta-out (LODO) --- coupled with area-of-applicability (AOA) masking
and distribution-free conformal prediction intervals. \PLACE{For SOC, a model
scoring $R^2\approx0.65$ under random cross-validation collapses to a median
LODO $R^2$ of $-1.1$, and every held-out delta falls entirely outside the AOA,
even with climate, geomorphic and tidal covariates included.} We show that the
failure is not explained by covariate-shift magnitude alone, implicating
delta-specific (concept-level) controls on carbon density. \PLACE{AGB and
total-carbon results.} Our results quantify a hard limit on the geographic
transferability of blue-carbon stock models and argue that crediting-grade
estimates for an unsampled delta require local calibration data.
\end{abstract}

\vspace{0.5em}
\noindent\textbf{Keywords:} blue carbon; mangroves; soil organic carbon;
out-of-distribution generalization; area of applicability; conformal prediction;
spatial cross-validation.

% ---------------------------------------------------------------------
\section{Introduction}
\PLACE{Open with the stakes: mangrove carbon density, the SOC-dominated stock
($\sim$10 of 11.7~Pg C belowground), and the rise of global wall-to-wall maps
(Sanderman 2018 SOC 30~m; Simard 2019 AGB 30~m) now feeding carbon markets and
Bangladesh Delta Plan-type policy.}

\PLACE{Problem: these maps report in-distribution skill. Spatial autocorrelation
inflates apparent accuracy \citep{ploton2020}; nobody has asked the
out-of-distribution question --- can blue-carbon knowledge from sampled deltas
predict an unsampled one? This is exactly the question a carbon buyer faces.}

\PLACE{Gap and contribution (enumerate):}
\begin{enumerate}
  \item A reproducible \emph{multi-delta benchmark} with frozen leave-one-delta-out
        splits for mangrove total ecosystem carbon (AGB-C + SOC).
  \item Three-tier validation that exposes the spatial-autocorrelation inflation
        and quantifies the \emph{transfer gap}.
  \item Trustworthiness layers --- area of applicability \citep{meyer2021} and
        conformal intervals --- that tell a practitioner \emph{where} a global
        model may be used.
  \item A mechanistic diagnostic separating covariate shift from concept shift.
\end{enumerate}

% ---------------------------------------------------------------------
\section{Results}

\subsection{A multi-delta blue-carbon benchmark}
\PLACE{Describe the frozen registry: candidate deltas, GMW-clipped mangrove area,
usable in-situ label counts, the core/prediction-only/excluded roles, continental
spread.} The v1 SOC benchmark comprises eight core deltas spanning four continents
(Table~\ref{tab:deltas}, Fig.~\ref{fig:map}); median 0--100~cm SOC stocks range from
\SI{61}{\mega\gram\per\hectare} (Zambezi) to \SI{519}{\mega\gram\per\hectare}
(Musi--Banyuasin) despite near-uniform tropical climate, signalling strong
delta-specific control.

\begin{table}[t]
\centering
\caption{Core (leave-one-delta-out) deltas in the v1 SOC benchmark: mangrove area
(GMW v3, 2020, equal-area), usable in-situ SOC cores, and median 0--100~cm SOC stock
(Mg\,ha$^{-1}$). Auto-generated from \texttt{data/processed/delta\_registry.csv}.}
\label{tab:deltas}
\begin{tabular}{llrrr}
\toprule
Delta & Continent & Area (km$^2$) & SOC cores & Median SOC \\
\midrule
Sundarbans & Asia & 6052 & 11 & 200 \\
Amazon Amapa & S.America & 4462 & 69 & 141 \\
Everglades & N.America & 1845 & 7 & 309 \\
Musi Banyuasin & Asia & 1637 & 37 & 519 \\
Saloum Gambia & Africa & 1066 & 36 & 252 \\
Mekong & Asia & 1021 & 119 & 270 \\
Zambezi & Africa & 847 & 12 & 61 \\
Rufiji & Africa & 412 & 96 & 158 \\
\bottomrule
\end{tabular}

\end{table}

\begin{figure}[t]
\centering
\includegraphics[width=\linewidth]{fig_map.pdf}
\caption{The eight core leave-one-delta-out deltas. Marker area scales with mangrove
extent (GMW v3, 2020); colour shows median in-situ SOC$_{0\text{--}100}$. Spread across
five continents with strongly varying soil-carbon density.}
\label{fig:map}
\end{figure}

\subsection{Standard validation overstates transferable skill}
A gradient-boosting model that explains \PLACE{65}\% of SOC variance under random
$k$-fold cross-validation drops to slightly negative skill under spatial-block CV and
collapses to strongly negative skill under leave-one-delta-out (Table~\ref{tab:tiers},
Fig.~\ref{fig:tiers}). The implied \emph{transfer gap} (random minus LODO) is
\PLACE{6.5} $R^2$ units --- the magnitude by which conventional validation overstates
out-of-distribution skill.

\begin{table}[t]
\centering
\caption{Three-tier validation of the SOC model. Random $k$-fold is inflated by
spatial autocorrelation \citep{ploton2020}; leave-one-delta-out (LODO) is the
out-of-distribution test. Auto-generated from \texttt{soc\_lodo\_results.json}.}
\label{tab:tiers}
\begin{tabular}{lrr}
\toprule
Validation tier & Ridge $R^2$ & Grad.\ boosting $R^2$ \\
\midrule
T1 random $k$-fold & +0.31 & +0.65 \\
T1 site-grouped $k$-fold & +0.10 & +0.16 \\
T2 spatial-block & -0.09 & -0.69 \\
T3 leave-one-delta-out (median) & -2.05 & -1.71 \\
\bottomrule
\end{tabular}

\end{table}

\begin{figure}[t]
\centering
\includegraphics[width=0.62\linewidth]{fig_tiers.pdf}
\caption{Model skill ($R^2$) under the three validation tiers. The drop from random
$k$-fold to leave-one-delta-out quantifies the inflation of conventional validation.}
\label{fig:tiers}
\end{figure}

\subsection{Every unsampled delta lies outside the area of applicability}
For all eight held-out deltas, \emph{zero} percent of the mangrove area falls inside
the model's area of applicability (Table~\ref{tab:perdelta}, Fig.~\ref{fig:aoa}): each
delta is a separated cluster in covariate space, so every cross-delta prediction is an
extrapolation. Conformal intervals targeting 90\% coverage realise \PLACE{0--90}\%
empirical coverage, confirming that distribution-free guarantees break under this
shift. This AOA result is model-agnostic and robust across all feature sets.

\begin{figure}[t]
\centering
\includegraphics[width=\linewidth]{fig_aoa.pdf}
\caption{Per-delta leave-one-delta-out error (left) and fraction of each delta inside
the area of applicability (right). Every delta lies entirely outside the AOA.}
\label{fig:aoa}
\end{figure}

\begin{table}[t]
\centering
\caption{Per-delta leave-one-delta-out results (gradient boosting, full covariate set):
sample size, $R^2$, RMSE (Mg\,ha$^{-1}$), fraction inside AOA, and empirical conformal
coverage (target 0.90). Auto-generated from \texttt{soc\_lodo\_results.json}.}
\label{tab:perdelta}
\begin{tabular}{lrrrrr}
\toprule
Delta & $n$ & $R^2$ & RMSE & AOA in & Conf.\ cov. \\
\midrule
Amazon Amapa & 69 & -0.08 & 55 & 0.00 & 0.90 \\
Everglades & 7 & -1.42 & 110 & 0.00 & 0.86 \\
Zambezi & 12 & -2.09 & 112 & 0.00 & 0.42 \\
Saloum Gambia & 36 & -1.30 & 116 & 0.00 & 0.69 \\
Rufiji & 96 & -2.01 & 117 & 0.00 & 0.71 \\
Sundarbans & 11 & -33.59 & 121 & 0.00 & 0.00 \\
Mekong & 119 & -1.11 & 144 & 0.00 & 0.60 \\
Musi Banyuasin & 37 & -4.83 & 205 & 0.00 & 0.86 \\
\bottomrule
\end{tabular}

\end{table}

\subsection{Adding geomorphic and tidal forcing does not restore transfer}
\PLACE{The climate $\to$ +geomorphic $\to$ +tidal (EOT20) progression; gap barely
moves. Argues the limit is not merely a missing-covariate artifact.}

\subsection{Covariate shift does not fully explain which deltas fail}
\PLACE{The diagnostic: weak/non-significant correlation between shift magnitude
(dissimilarity index, energy distance) and LODO error, with the small-$n$ caveat;
physically coherent per-delta shift axes (tidal for Amazon/Rufiji/Musi,
precipitation seasonality for Saloum, monsoon thermal regime for Sundarbans).
Points to concept shift. Reference Table in SI.}

\subsection{Aboveground and total ecosystem carbon}
\TODO{AGB labels (Simard 2019) acquired for all core deltas; AGB LODO models and
AGB-C + SOC $\to$ total ecosystem carbon with propagated uncertainty are the next
analysis. Section reserved.}

% ---------------------------------------------------------------------
\section{Discussion}
\PLACE{Interpretation: a hard transferability limit for blue-carbon stock models;
implications for carbon crediting (no crediting-grade estimate for an unsampled
delta without local cores); why concept shift (delta-specific sediment/salinity/
hydroperiod controls) is the likely mechanism; what would close the gap (local
calibration, hierarchical/partial pooling, transfer learning).}

\PLACE{Limitations: SOC label sparsity and depth harmonization; small per-delta
$n$; provisional AOA threshold; tidal from EOT20 rather than a process model;
v1 is SOC-led with AGB pending.}

\PLACE{Outlook: total-carbon mapping, the change/emissions and SLR-projection
modules, and the framework's portability to any delta via a bounding box.}

% ---------------------------------------------------------------------
\section{Methods}

\subsection{Study systems and the delta registry}
\PLACE{Pre-registered selection criteria (area, deltaic setting, label
availability, continental spread, non-overlap to prevent LODO leakage). GMW v3
2020 extent clipped per delta in an equal-area projection; usable-core gate from
the harmonized SOC labels.}

\subsection{In-situ carbon labels}
\PLACE{SOC: Coastal Carbon Network Data Library v1.6.0; mangrove cores integrated
to a standardized 0--100~cm stock (mass-preserving; bulk-density $\times$ carbon
fraction; OM$\to$C fallback ratio; extrapolation flags). AGB: Simard et al. 2019
30~m, AGB$\to$C $\times 0.47$.}
\begin{equation}
\mathrm{SOC}_{0\text{--}100} \;=\; 100 \sum_i \mathrm{BD}_i \, f_{C,i}\, \Delta z_i
\quad [\mathrm{Mg\,C\,ha^{-1}}],
\label{eq:soc}
\end{equation}
\PLACE{define terms, units, depth handling.}

\subsection{Covariates}
\PLACE{Climate (CHELSA bio1--19); soil context and SOC prior (SoilGrids, GSOCmap);
land cover (Copernicus); geomorphic position (distance-to-coast/-river, Natural
Earth); tidal range and form factor from the EOT20 ocean-tide model. External
driver lake referenced via \texttt{MDBC\_LAKE}.}

\subsection{Models and three-tier validation}
\PLACE{Ridge baseline and histogram gradient boosting on $\log(1+\mathrm{SOC})$.
T1 random $k$-fold; T2 spatial-block CV \citep{roberts2017}; T3 leave-one-delta-out.
Transfer gap $=$ T1$-$T3.}

\subsection{Area of applicability}
\PLACE{Importance-weighted dissimilarity index to the training feature space with
an outlier-aware threshold \citep{meyer2021}; report fraction of each held-out
delta inside the AOA.}

\subsection{Conformal prediction intervals}
\PLACE{Split-conformal on held-out calibration residuals for finite-sample marginal
coverage; report empirical coverage per delta.}

\subsection{Transfer diagnostic}
\PLACE{Per-delta covariate shift via mean dissimilarity index and multivariate
energy distance; Spearman correlation against LODO error; standardized
mean-difference ranking to name the dominant shift axis.}

% ---------------------------------------------------------------------
\section*{Data availability}
\PLACE{All source datasets are public (CCN, GMW v3, Simard 2019, EOT20, CHELSA,
Natural Earth); per-source provenance manifests and derived label/feature tables
are released with the code.}

\section*{Code availability}
\PLACE{Reproducible pipeline (Stage 0--2), pinned environment and a single-command
runner are available at \PLACE{repository URL on publication}.}

\section*{Acknowledgements}
\PLACE{Funding and data-provider acknowledgements.}

\section*{Author contributions}
\PLACE{M.R.H.: conceptualization, methodology, software, analysis, writing.}

\section*{Competing interests}
The author declares no competing interests.

% ---------------------------------------------------------------------
\bibliographystyle{plainnat}
\bibliography{references}

\end{document}
